\honest{Semantic Stopsigns}{Eliezer Yudkowsky}{2007}

A child asks where a rock came from. From the boulder; the boulder came from the mountain; the mountain is the bones of the giant Ymir; Ymir came from the abyss Ginnungagap. And where did Ginnungagap come from? ``Never ask that question.''

Science traces events back to the Big Bang, but the physical laws still seem to call for an explanation; if we are a computer simulation, the computer runs on some other world's laws, and where did those come from? At this point some people say ``God!'' I ask what could make anyone think this even helped, since the next question is ``Where did God come from?'' Saying that God is uncaused, or created Himself, leaves us just where ``Time began with the Big Bang'' did: we can still ask why the whole system exists, or ``why some events but not others are allowed to be uncaused.'' \nb{Defenders of the cosmological argument answer that question: only what is contingent, or what begins to exist, needs a cause. The reply is disputed, and the post does not mention it.}

I am not discussing the First Cause itself, only asking why anyone thinks ``God!'' resolves it. My answer is that saying ``God!'' ``is a way of belonging to a tribe,'' so people say it as often as they can, even about why a hurricane struck New Orleans. I give no evidence for this. But even if God exists, saying so leaves the First Cause as puzzling as before; how could anyone not notice?

Jonathan Wallace suggested that ``God!'' is a semantic stopsign: not a claim, but a signal not to think past this point. \nb{Wallace used the phrase ``semantic stop-sign'' in an essay written after the September 11 attacks, and made the same point in 1995 with a similar chain of questions about a made-up creator called ``Og.''}

Of course you, ``being a good and proper atheist,'' would never do that. But there are other stopsigns. New technologies raise hard policy questions: should a government supervise parents' choice of their child's genes, could parents choose genes for schizophrenia, should governments pay for intelligence enhancement to prevent a cognitive elite? For any proposed institution, charities or product-liability lawsuits for example, the next question is whether it will work. I know someone whose answer to all of them is ``Liberal democracy!'' ``That's it. That's his answer.'' Anyone who asks how democracies have done on such problems is, to him, an autocrat. Had someone proposed the Coca-Cola corporation instead, he would have asked all the obvious questions. \nb{That is what the author expects he would have asked. The post reports nothing he said about it.} In the same way, if you would question a claimed Mexican-American plot to remove Earth's oxygen but not a claimed corporate one, then ``Corporations!'' is a stopsign for you.

A caution: do not turn ``it's just a stopsign'' into a new all-purpose counterargument. No word is a stopsign in itself; it depends on the person, and strong feelings about something do not make it one. The word ``intelligence'' once was one for me. What marks a stopsign is ``failure to consider the obvious next question.'' \nb{By this test, a theist who gives the reply described in the earlier note has considered the obvious next question, and so is not using a stopsign, though the post offered ``God!'' as ``the obvious first example.''}
